\section{Preliminaries}

This section establishes foundational concepts and terminology required to understand the CAEM methodology, mathematical formulations, and experimental evaluation presented in subsequent chapters.

\subsection{Large Language Models and Transformer Architecture}

Neural networks that have been trained on quite big text collections to predict and produce natural language are known as large language models. The Transformer architecture \cite{NIPS2017_3f5ee243}, which processes sequences using a self-attention mechanism allowing all the tokens in the input to take account of every other token, is adopted by new LLMs. The encoder and decoder layers in this architecture are arranged in layers and each one has feed forward networks and multi-head self-attention.

In this study, we utilize Flan T5 \cite{JMLR:v25:23-0870}, an encoder-decoder model that integrates instruction based fine tuning and builds the existing T5 architecture. Encoder decoder models generate text step by step through the decoder after creating a bidirectional representation of the input through the encoder, in contrast to decoder only models like GPT that read text from left to right. Because the model is capable of fully understanding the question before generating an answer, this method is particularly effective for tasks involving answering questions.

Fine tuning adapts pre trained models to specific tasks or domains by continuing training on data specific to the task while using smaller learning rates. The model parameters $\theta$ are adjusted to reduce loss on the new data while still preserving the general abilities learned during pre training. Standard fine tuning updates all parameters, while parameter efficient approaches such as LoRA modify only a small portion of them.

\subsection{Hallucination Taxonomy}

Hallucinations in language model outputs can be understood along two main dimensions. The first dimension separates \textit{intrinsic hallucinations}, which directly contradict information provided in the source input, from \textit{extrinsic hallucinations}, which introduce information that goes beyond what the source provides, regardless of whether it also happens to conflict with the source \cite{ji2023survey}. For instance, if a biography states that ``Einstein was born in Germany'' and the model generates ``Einstein was born in France,'' this is an intrinsic hallucination because it contradicts the stated source. But if the model generates ``Einstein played violin exceptionally well'' without any evidence in the source, this is an extrinsic hallucination even if the claim happens to be true --- the defining characteristic is that it cannot be verified from the input, not that it contradicts anything in it.

The second dimension distinguishes \textit{factuality hallucinations}, which conflict with verifiable world knowledge, from \textit{faithfulness hallucinations}, which contradict user input or task context \cite{10.1145/3649506}. If a model claims "photosynthesis produces oxygen from carbon dioxide and water" and demonstrates factual accuracy. But on the other hand it it claims "photosynthesis occurs in mitochondria" then this demonstrates factual hallucination. Faithfulness focuses on whether the output of the model remains consistent with the supplied input or context, especially when summarizing documents or following instructions.

This research focuses on factuality hallucinations which occurs in knowledge based reasoning tasks. Here correct answers can be verified against trusted sources or logical rules and regulations. But faithfulness to fetch evidences also matter for etrieval augmented generation. Here outputs must be grounded in retrieved documents.

\subsection{Natural Language Inference}

Natural Language Inference (NLI) is the task of determining whether a hypothesis sentence logically follows from a premise sentence \cite{10.1007/11736790_9}. Given premise $P$ and hypothesis $H$, an NLI model classifies the relationship as:
\begin{itemize}
    \item \textbf{ENTAILMENT}: $P$ implies $H$ is true
    \item \textbf{CONTRADICTION}: $P$ implies $H$ is false  
    \item \textbf{NEUTRAL}: $P$ neither confirms nor denies $H$
\end{itemize}

For example, if we assume "The Beatles were formed in Liverpool in 1960" and hypothesis is "The Beatles were an English band," then this relationship is ENTAILMENT. Given hypothesis "The Beatles were formed in London," the relationship is CONTRADICTION. But if the given hypothesis is "The Beatles recorded many albums," then the relationship is NEUTRAL.

We employ RoBERTa Large MNLI \cite{liu2019roberta}, which is a model that is fine tuned on the Multi Genre Natural Language Inference dataset and it can achieve an accuracy of approximately 90.8\% on the MNLI matched development set. NLI serves two purposes in CAEM. Firstly, it verifies factual correctness by checking whether generated claims are rewuired by retrieved evidence, and secondly,  t helps group answers that mean the same thing by testing for bidirectional entailment. If P entails H and H entails P, then P and H are considered semantically equivalent.

\subsection{Sentence Embeddings and Semantic Similarity}

Sentence embeddings map the length of variable text sequences to fixed dimensional dense vectors in a semantic space where similar meanings have high cosine similarity. Sentence BERT (SBERT) \cite{reimers-gurevych-2019-sentence} produces encasings or embeddings by fine tuning BERT in a siamese network structure on semantic textual similarity tasks. The model we employ, all mpnet base v2, generates 768-dimensional embeddings.

Cosine similarity between vectors $\mathbf{v}_1$ and $\mathbf{v}_2$ is computed as:
\begin{equation}
\text{sim}(\mathbf{v}_1, \mathbf{v}_2) = \frac{\mathbf{v}_1 \cdot \mathbf{v}_2}{\|\mathbf{v}_1\| \|\mathbf{v}_2\|} = \cos(\theta)
\end{equation}
where $\theta$ is the angle between vectors. Similarity ranges from -1 (opposite) to 1 (identical), with 0 indicating orthogonality.

Semantic similarity is different from string similarity. For instance, the questions "Is 17 prime?" and "Is the number 17 a prime number?" have low string overlap but high semantic similarity (approximately 0.98). CAEM uses semantic similarity for key operations like memory retrieval (finding similar questions), novelty detection (avoiding redundant storage), and self-consistency verification (measuring reasoning chain agreement). The system uses FAISS (Facebook AI Similarity Search) \cite{8733051} for fast and efficient approximate nearest neighbor search, allowing it to retieve  sub millisecond findings over thousands of stored episodes.

\subsection{Self-Consistency Decoding}

Self consistency \cite{wang2023selfconsistency} improves reasoning quality by first generating multiple independent reasoning chains for the same question, then it selects the most common answer. Unlike standard greedy decoding which produces a single output, self-consistency samples $M$ chains using temperature $T > 0$ to introduce diversity.

For CAEM, we adapt self consistency technique for verification rather than choosing an answer. Given question $q$ and generated answer $a$, we generate $M$ additional chains and measure agreement via average pairwise similarity:
\begin{equation}
u_{\text{SC}} = \frac{1}{M^2} \sum_{i=1}^{M} \sum_{j=1}^{M} \text{sim}(\mathbf{v}_i, \mathbf{v}_j)
\end{equation}
where $\mathbf{v}_i$ is the embedding of chain $i$. High consistency ($u_{\text{SC}} > 0.85$) indicates robust reasoning where multiple paths converge, while low consistency suggests uncertainty or errors.

\subsection{Catastrophic Forgetting and Continual Learning}

Catastrophic forgetting \cite{MCCLOSKEY1989109} occurs when neural networks learn new tasks and suddenly lose their ability to perform tasks on previously learned materials. Standard gradient descent focuses only on improving performance for the current task and does not protect earlier knowledge. This causes the new parameter to update that help the new task but weakens old tasks. Fine tuning LLMs on domain specific data can reduce accuracy on general benchmarks by 30-50\% without mitigation strategies.

Elastic Weight Consolidation (EWC) \cite{doi:10.1073/pnas.1611835114} addresses this by identifying essential parameters for previous tasks using the Fisher Information Matrix, then adding quadratic penalties preventing large updates to important parameters:
\begin{equation}
\mathcal{L} = \mathcal{L}_{\text{new}} + \frac{\lambda}{2} \sum_{i} F_i (\theta_i - \theta_i^*)^2
\end{equation}
where $F_i$ measures parameter importance, $\theta^*$ are previous parameters, and $\lambda$ controls regularization strength. Simpler methods include applying L2 regularization that can pull the model back toward its original parameters and mixing new training data with samples from earlier tasks. CAEM uses these mitigation strategies to maintain broad language abilities while fine tuning on verified examples. This addresses the challenge of learning from imperfectly verified data (85-90\% accuracy) without degrading through error reinforcement.

The conceptual basis for understanding CAEM's multilayer verification, estimation of confidence, memory architecture and self improvement structure are discussed and provided in these groundworks. Following sections describes ow prior research develops and validates these individual components, preparing for their integration in Chapter 4.

\section{Review of Existing Research}

This section surveys prior work across six areas essential to understanding CAEM's design: hallucination characterization and measurement, verification methods for reasoning validation, confidence estimation techniques, memory-augmented architectures, continual learning approaches, and evaluation benchmarks. Each subsection identifies key contributions and limitations that motivate CAEM's architectural choices.

\subsection{Hallucination Problem and Measurement}

Understanding the nature and characteristics of hallucination in LLMs requires methodologies like taxonomic frameworks and quantitative measurement. Early studies on hallucination shows summarization and data to text generation. Ji et al. \cite{ji2023survey} published the first major survey in ACM Computing Surveys. It establishes a study differentiating intrinsic hallucinations that contradicts source material from extrinsic hallucinations. Even though the claims were unverified. Their work shows that hallucinations occurs at several stages. It includes data collection that embeds factual errors, training procedures that reinforce misleading correlations, and inference methods that fail to account for uncertainty.

Huang et al. \cite{10.1145/3649506} extended this study specifically for large language models in ACM Transactions on Information Systems. It is extended by adding the differencs between factuality and faithfulness. Their analysis of causes across the model lifecycle suggests various groundbreaking research. It suggests that effective mitigation should have intervention at several stages instead of relying solely on post hoc detection. This directly supports CAEM's multi component strategy which combines memory, verification, and iterative learning.

Quantitative measurement studies of hallucination prevalence reveals concerning statistics across domains and model scales. Lin et al. \cite{lin-etal-2022-truthfulqa} introduced TruthfulQA at ACL2022, a benchmark of 817 questions across 38 categories which has been specifically designed to elicit imitative falsehoods where models reproduce common human misconceptions found in training data. Their result demonstrates an inverse scaling effect: GPT-3 175B achieved only 21-25\% on the truthful and informative metric (or 58\% on truthfulness alone) compared to 94\% human performance. Larger models often performed worse than smaller ones. This proves that only scaling is not sufficient for truthfulness and thus it needs strong motivation for architectural innovation.

The inverse scaling result has major result for hallucination mitigation strategies. If adding parameters reduces truthfulness on some tasks then mitigation strategies must operate beyond parameter growth. CAEM addresses this by incorporating external verification and memory components that operate independently of model scale, providing corrections and verified knowledge regardless of base model size.

Min et al. \cite{min-etal-2023-factscore} introduced FactScore at EMNLP 2023. It provides smooth evaluation through the process of breaking down complex. Instead of assigning a single label to check correctness to an entire answer, FactScore computes the percentage of claims that is supported by sources. As a result, this method achieves less than 2\% error rate compared to human evaluation. Applied to biography generation, ChatGPT achieved only 58\% FActScore, which indicates that even the highly capable models produce substantial proportions of unsupported atomic facts. This findings highlights the need for fact level verification rather than holistic response evaluation.

Domain specific studies show even higher hallucination rates in specefic contexts. Alkaissi and McFarlane \cite{alkaissi2023artificial} studied by making ChatGPT response to 25 clinical questions and found 69-77\% contained hallucinated medical information. Cheli et al. \cite{info:doi/10.2196/53164} studied hallucination in medical literature queries. He reported rates of 69.6\% for GPT-3.5,  28.6\% for GPT-4, and 91.4\% for Google Bard.

Recent benchmark provided comprehensive evaluation across Several dimensions. One of them being legal sector. Factual accuracy has been measured in legal sector. According to Google's FACTS benchmark (2025), it shows that even top performing systems maintain 6.4\% hallucination rates on legal queries. Furthermore, the Vectara Hallucination Leaderboard tracks document summarization hallucination rates ranging from 3\% for GPT-4 to 27.2\% for Google PaLM 2 Chat. It thus demonstrates wide variation among models and confirm that hallucination remains an issue even in state of the art systems.

The information above has revealed three critical takeaways that shape CAEM's designe. Firstly, hallucination occurs across models, tasks, and domains. Depending on context the rates ranging from 3\% to over 90\% . Secondly, larger models do not automatically produce fewer hallucinations. This implies that architectural intervention is necessary. Lastly, verifying information at the level of atomic facts is far more informative than evaluating full responses, motivating CAEM's commitment to fact level NLI based verification.

\subsection{Verification Methods}

Verification methods are used to check whether a model's reasoning is correct before the information is stored or shown to users. This subsection explains four approaches that form the basis for CAEM's multi layer verification architecture.

Chain of thought prompting, which was introduced by Wei et al. \cite{NEURIPS2022_9d560961} at NeurIPS 2022, improves reasoning quality by prompting models to generate the next steps before final answers. Showing these intermediate steps leads to stronger reasoning and provides better starting points for verification. When tested on the PaLM 540B model, chain of thought greatly improved performance on math reasoning tasks, increasing accuracy on GSM8K from about eighteen percent to over fifty eight percent. The authors found that this technique showed an emergent ability threshold around 100 billion parameters for \textit{base models} trained on raw text, below which chain-of-thought provided minimal benefit. However, this threshold does not apply to instruction-tuned models, which learn to follow step-by-step reasoning patterns during the instruction-tuning phase itself. Chung et al.\ \cite{JMLR:v25:23-0870} demonstrated that Flan-T5 models show substantial chain-of-thought benefits at 780 million parameters after instruction tuning. Since CAEM uses Flan-T5-Large (780M parameters), which is an instruction-tuned model, chain-of-thought prompting is effective despite being well below the 100B base-model threshold.

For CAEM, chain of thought serves two purposes. Firstly, it improves base model reasoning quality. It also provides better starting points for verification. Secondly, it gives explicit reasoning chains that enables step by step verification through NLI which also checks whether each reasoning step is supported by evidence. The explainability of chain of thought reasoning proves essential for diagnosis when verification fails, enabling error analysis that identifies which reasoning steps caused failures.

Wang et al. \cite{wang2023selfconsistency} expanded this idea in 2023 by introducing self consistency. Instead of producing a single reasoning path, the model generates several different explanations and then selects the most common answer. The key idea is that difficult and complex problems usually has multiple valid reasoning paths which leads to the same correct answer. On the other hand, incorrect answers rarely achieve consistency across independent derivations. Quantitative improvements with PaLM 540B demonstrate substantial gains: +17.9\% on GSM8K, +11.0\% on SVAMP, +12.2\% on AQuA, and +6.4\% on StrategyQA.

CAEM uses self consistency in different ways. It chooses reasoning chain agreement through semantic similarity rather than selecting the majority answer. If the chains mostly agree, the reasoning is then considered reliable and suitable for storage. But if they differin opinion then the system treats the result as uncertain and may apply further checks or reject it. As a result, this method is efficient since it only requires a small number of extra generations which makes it practical for regular verification. And it also provides useful confidence signals.

Another important method is semantic entropy introduced bt Farquhar et al. \cite{farquhar2024semantic} in 2024. It represents the most rigorous treatment of uncertainty based hallucination detection. Their key innovation addresses a primary limitation of token level entropy: Answers that express the same idea in different wording are grouped together using inference checks, and uncertainty is measured across these meaning based groups.

The methodology achieves AUROC of 0.790 for frabrication detection on question answering tasks. It substantially outperforms token entropy (0.691) and verbalized probability methods (0.698). Even though the performance remains robust across model families including LLaMA 2 (7B-70B), Falcon (7B-40B), and Mistral 7B, the Nature publication validates semantic entropy as a cutting edge technique for uncertainty quantification, motivating its inclusion in CAEM despite computational overhead.

Semantic entropy serves as the third verification layer for CAEM. It catches reports where both NLI and self consistency fail. Let us consider a model which is asked "What is the capital of Australia?" that consistently always generates "Sydney" across multiple chains with high individual probabilities. NLI verification may pass if the model generates factually correct supporting statements about Sydney. Here, self consistency shows perfect agreement. But still the semantic entropy will generate 10 diverse samples and will measure cluster entropy. This may reveal that the model also produces "Canberra" and "Melbourne" at high temperature thus indicating underlying uncertainty. This detection capability for systematic confabulation justifies the computational cost.

Natural language inference provides direct verification by checking whether a claim is supported by evidence. Thorne et al. \cite{thorne-etal-2018-fever} introduced the FEVER dataset at NAACL 2018 which formulates fact verification as NLI with 185,445 claims generated from Wikipedia. Systems must classify claim evidence relationships as SUPPORTS, REFUTES, or NOT ENOUGH INFO. This three way classification would map directly to CAEM's storage decisions. The decision would be to support claims merit storage, refute claims are rejected, and insufficient evidence trigger uncertainty handling.

Pre trained NLI models such as RoBERTa-Large-MNLI has achieved approximately 90\% accuracy on the MNLI benchmark. It thus provides reliable verification for factual claims. But the limitation is that even though NLI verifies individual factual statements, they may miss logical reasoning errors that do not contradict facts. For example, a model could state "Sydney is Australia's largest city" which is true and conclude "therefore Sydney is the capital" which is an invalid reference. NLI would verify the factual premise while missing the logical error thus motivating the need for complementary verification layers.

The integration of these four methods provides complementary strengths. Chain of thought improves reasoning interpretability. On top of that, self consistency catches inconsistent logic. Furthermore, Semantic entropy detects confabulation and NLI grounds factual claims. Together, they achieve the 85-90\% verification accuracy target that enables reliable quality control for episodic memory, addressing different failure modes that individual methods miss.

\subsection{Confidence Estimation}

Reliable confidence estimation helps systems tell the difference between answers that can be trusted and ones that need extra checking from outside sources. This section looks at basic methods for measuring uncertainty and explains how these ideas have been adapted for use with large language models.

Gal and Ghahramani \cite{pmlr-v48-gal16} laid the theoretical foundations for practical uncertainty estimation in their ICML 2016 paper "Dropout as a Bayesian Approximation." Their key contribution explains that dropout usually during training can be understood as approximate Bayesian inference over model parameters. By keeping dropout active during testing and running the model multiple times, it becomes possible to measure how much the outputs vary.This variation can be used to estimate epistemic uncertainty, which refers to uncertainty caused by limited knowledge about the model itself. And this is different from aleatoric uncertainty, which comes from noise or randomness in the data. But it cannot be reduced.

The theoretical framework shows that a network with dropout before each weight layer approximates a deep Gaussian process. This predicted average and variance can be approximated through Monte Carlo sampling. It describes as performing $K$ forward passes with different dropout masks and computing output variance. For better understanding, this provides uncertainty estimates without requiring group training or modifying the loss function. Because this approach works with any network that already uses dropout, it is especially useful for adapting existing large language models.

For CAEM, MC Dropout serves as the second confidence signal. It quantifies model uncertainty through output variance. Here, high variance indicates the model is uncertain about its output and suggests that verification is needed. On the other hand, low variance indicates confident generation. Even though this alone is insufficient as models can be confidently wrong, the computational cost of 5-10 forward passes is non trivial but acceptable for Tier 2 and Tier 3 routing where generation already incurs substantial cost.

A major problem with neural network confidence is that models are often too confident in their answers. Guo et al. \cite{pmlr-v70-guo17a} demonstrated this at ICML 2017. They showed modern deep networks predict confidence scores substantially which exceeds actual accuracy rates. For example, predictions at 80\% confidence may achieve only 60\% accuracy and create a  dangerous miscalibration for decision making. They introduced Expected Calibration Error (ECE) as the standard metric, dividing predictions into bins by confidence and measuring the gap between confidence and accuracy within each bin.

Their proposed solution, temperature scaling, achieves remarkable effectiveness with minimal complexity. The method divides logits by a learned scalar temperature $T > 0$ optimized on a validation set: $\hat{y} = \sigma(\mathbf{z}/T)$ where $\mathbf{z}$ are logits and $\sigma$ is softmax. Temperature scaling preserves accuracy (argmax unchanged) of classification while calibrating probabilities. Even though only a single parameter is used, it outperforms more complex methods including histogram binning, isotonic regression, and Bayesian approaches. The simplicity and effectiveness make temperature scaling the standard post hoc calibration method.

CAEM uses temperature scaling as the final step in estimating confidence. This step adjusts the combined confidence score from all four signals so that it matches real accuracy observed on validation data. As a result, when the system reports eighty percent confidence, about eighty percent of those predictions are actually correct. This makes it possible to set reliable thresholds for deciding when information can be trusted and when additional verification is needed.

Kadavath et al. \cite{kadavath2022language} from Anthropic explored self evaluation in large language models at NeurIPS 2022. They introduced two methods for extracting confidence from LLMs themselves. First is P(True) asks models to evaluate the probability their own previous answers are correct. And, P(IK) or "I Know" prompts models to predict confidence before generating answers. Their findings reveal nuanced calibration properties. Larger models tend to show better calibration on multiple choice questions which means their confidence levels line up more closely with how often they are actually correct. However, reinforcement learning from human feedback can make calibration worse, since it often rewards answers that sound helpful or confident rather than ones that are strictly accurate.

The idea that reinforcement learning from human feedback can hurt calibration is important. Models are trained to be helpful. But they may express high confidence to appear authoritative even when uncertain which is exactly the opposite of desired behavior for high stakes applications. Because of this conflict between sounding helpful and being truthful, CAEM relies on several outside confidence signals instead of trusting the model's own confidence alone.

Xiong et al. \cite{xiong2023can} conducted a detailed study of confidence estimation methods for large language models. They provided comprehensive approaches into white box which use internal information like token logits and sequence probabilities, and black box methods, which rely on external signals such as verbalized confidence or consistency checks. Their systematic evaluation across multiple benchmarks reveals striking findings. That is, even the strongest single methods only reached about 0.50 to 0.60 AUROC for telling correct answers from incorrect ones, which is just slightly better than random guessing.

CAEM's multi signal approach is influenced by limited discriminative ability with individual signals. The system combines token probability, MC Dropout, self consistency, and semantic entropy to take advantage of different types of information. Token probability reflects model's direct confidence, MC Dropout measures parameter uncertainty. It also captures self consistency reveals reasoning stability, and semantic entropy detects meaning-level confusion. The weighted combination achieves the 89\% routing accuracy necessary for reliable tier selection.

Recent work by Tian et al. \cite{tian2023just} questions if models genuinely know when they don't know. It also questions if the models only match patterns for uncertain expressions from training data.Their study found that verbalized confidence often reflects patterns in the training data rather than the model's actual knowledge or uncertainty. This reinforces CAEM's strategy of combining implicit signals like token probability, dropout, and self consistency with explicit methods, instead of relying only on verbalized confidence, which can be misleading or superficial.

The confidence estimation literature establishes several key findings for CAEM's design. Firstly, relying on a single confidence signal does not provide enough accuracy for safe routing. It in return makes combining multiple signals necessary. Secondly, systematic overconfidence requires post hoc calibration through temperature scaling. Thirdly, white box methods that access internal model information offer complementary insights compared to black box approaches which is based on consistency. Lastly, RLHF training may degrade calibration, pre-RLHF models or instruction-tuned models like Flan T5 may give more reliable confidence signals. Together, these points shape CAEM’s four signal confidence estimation system and its temperature scaled calibration process.

\subsection{Memory-Augmented Architectures}

External memory architectures give neural networks a place to store information outside their usual parameters. It  helps them learn and fetch facts. This subsection explains how these ideas developed. Starting with early systems that used differentiable memory and ending with the modern approach known as retrieval augmented generation.

Graves et al. \cite{graves2014neural} introduced Neural Turing Machines (NTMs) in 2014. They ventures into the idea of giving neural network external memory. And as a result, they could read from and write to in a smooth and trainable way. The architecture combines a controller network with an addressable memory matrix. This supports both content based addressing thats is done via cosine similarity and location based addressing, done via shift operations. The key innovations include read and write operations. It will enable end to end training via backpropagation. As a result, NTMs successfully learned algorithmic tasks. These includes copying, sorting, and associative recall.

The importance of CAEM is that it shows how external memory can give a model abilities that normal models do not have. NTMs were mainly used for learning algorithms. But, it can store facts as well. CAEM uses this idea in its episodic memory. It stores verified question and answer pairs along with the reasoning behind them. This lets the model call them up when needed without any delay.

Weston et al. \cite{weston2015memory} introduced Memory Networks at ICLR 2015. Their model scaled to millions of memory slots. It also showed that a system could perform several rounds of reasoning to answer a question. The architecture consists of four learned components. The components areL: turning the input into features, updating the memory, turning the memory output into features, and finally producing a response. Furthermore, Memory Networks performed strongly on bAbI question. With answering tasks that required reasoning over multiple facts. Sukhbaatar et al. \cite{NIPS2015_8fb21ee7} extended this  model with end to end Memory Networks at NeurIPS 2015. He removed the need for supervision on retrieval. He did this model by continuous attention mechanisms. This described how CAEM builds its own style of step by step reasoning using stored knowledge.

Graves et al. \cite{graves2016differentiable} presented the Differentiable Neural Computer (DNC) in Nature 2016. It built on the idea of Neural Turing Machines but added smarter memory tools like dynamic allocation. As a result it keep track of the order in which things were written. DNCs achieved 3.8\% average error on bAbI tasks compared to what it achieved for memory networks which is 6.4\%. It clearly shows that calculated memory management improves performance.

Lewis et al. \cite{NEURIPS2020_6b493230} introduced Retrieval-Augmented Generation (RAG) at NeurIPS 2020. They showed that combining what a model knows in its parameters with information stored outside the model can make its answers much more accurate. RAG adds a large searchable index of Wikipedia to a BART model with four hundred million parameters. It uses Dense Passage Retrieval to look up relevant documents. Then it works out the final answer by considering the generation probability across all the documents it retrieved. The architecture achieves substantial improvements over parametric-only models: 44.5\% exact match on Natural Questions (versus 41.5\% prior SOTA), 68.0\% on TriviaQA (versus 60.5\%), and 45.5\% on WebQuestions (versus 44.7\%).

Critically, RAG demonstrates index hot-swapping: updating the document index without retraining achieves 70\% accuracy on questions about updated world knowledge, proving that non-parametric memory can be updated without expensive retraining. This property directly motivates CAEM's episodic memory, which can be updated continuously through verification without model fine-tuning, though CAEM additionally employs iterative fine-tuning to consolidate learned patterns.

The limitation of standard RAG for CAEM's purposes is that retrieved knowledge is not verified before use and does not accumulate verified reasoning chains. CAEM extends RAG by adding verification layers and storing verified outputs for future retrieval, creating a growing knowledge base rather than relying solely on static external corpora.

Johnson et al. \cite{8733051} developed FAISS o handle similarity search at the scale of billions of item in 2019. It resulted in efficient implementation of memory augmented systems. The library achieves 8.5x speedup over prior GPU methods through optimized k-selection algorithms. Furthermore,FAISS also includes several types of indexes, such as IVF with product quantization, which allows the system to search in less than linear time. For CAEM, FAISS is the tool that makes real time memory search possible across about twenty thousand stored episodes. Keeping the search time below one tenth of a second is important for making the system fast enough to use in practice.

Reimers and Gurevych \cite{reimers-gurevych-2019-sentence} introduced Sentence BERT at EMNLP 2019. This model produced semantically meaningful sentence. This is achieved through embeddings that follows siamese network training. TThe speed improvement over regular BERT is huge. With standard BERT, finding the most similar sentence pairs in a group of ten thousand sentences takes about sixty five hours. With Sentence BERT, the same task takes only about five seconds because you can compare the embeddings directly instead of running a heavy cross encoder.This big jump in efficiency is what makes CAEMs memory search possible. Each query is turned into a 768-dimensional embedding, and the system can look through the entire memory index in real time.

These memory architecture developments establish the feasibility and benefits of external knowledge storage. NTMs prove differentiable memory enables structured reasoning. Memory Networks scale to millions of facts with multi-hop reasoning. RAG demonstrates that non-parametric knowledge improves factual accuracy and supports dynamic updates. FAISS and Sentence-BERT provide the efficient infrastructure for practical deployment. CAEM synthesizes these advances, employing verified episodic memory with FAISS indexing and Sentence-BERT embeddings, while extending beyond static retrieval to include verification, confidence-gated routing, and iterative improvement through fine-tuning on accumulated verified knowledge.

\subsection{Self-Improvement and Continual Learning}

Self-improvement mechanisms enable models to learn from interaction and feedback, while continual learning addresses the challenge of acquiring new knowledge without forgetting previous capabilities. This subsection reviews methods for iterative refinement and catastrophic forgetting mitigation that inform CAEM's self-improvement loop.

Kirkpatrick et al. \cite{doi:10.1073/pnas.1611835114} introduced Elastic Weight Consolidation (EWC) at PNAS 2017, providing a principled solution to catastrophic forgetting in neural networks. When learning sequential tasks, standard gradient descent optimizes exclusively for current task performance, causing parameter updates that improve new task accuracy while dramatically degrading previous task performance. EWC addresses this by estimating parameter importance for previous tasks using the Fisher Information Matrix, then adding quadratic penalties preventing large updates to important parameters.

The EWC loss function augments standard task loss with a regularization term:
\begin{equation}
\mathcal{L}_{\text{EWC}} = \mathcal{L}_{\text{task}} + \frac{\lambda}{2} \sum_{i} F_i (\theta_i - \theta_i^*)^2
\end{equation}
where $F_i$ approximates the Fisher Information measuring how much loss increases when parameter $\theta_i$ changes, $\theta^*$ are previous task parameters, and $\lambda$ controls regularization strength. Parameters important for previous tasks (high $F_i$) receive strong penalties against change, while less important parameters adapt freely to new tasks.

Applied to Atari game learning, EWC maintained approximately 90\% of single-task performance when sequentially learning 10 games, compared to severe degradation with standard SGD. For CAEM, EWC provides theoretical justification for L2 regularization toward base model parameters during fine-tuning. While full Fisher Information Matrix computation proves expensive for large models, simple L2 regularization $\frac{\lambda}{2}\|\theta - \theta_{\text{base}}\|^2$ approximates EWC's protective effect, preserving general capabilities while adapting to verified examples.

The practical implementation in CAEM employs $\lambda = 0.01$ based on continual learning literature, combined with mixed training data (90\% verified examples, 10\% general tasks) to maintain diverse capability. Early stopping based on MMLU score monitoring provides additional safety: if general capability retention drops below 93\%, training halts before catastrophic degradation occurs.

Natarajan et al. \cite{natarajan2013learning} provided theoretical foundations for learning with noisy labels at NeurIPS 2013. Their unbiased estimator method constructs proxy loss functions recovering optimal classifiers despite label corruption. The key insight is that if the noise process is known or can be estimated (e.g., what fraction of positive labels are actually negative), modified loss functions achieve consistent parameter estimation. Their method achieves more than 88\% accuracy with 40\% label corruption which substantially outperforms standard training. 

For CAEM, this framework directly applies to verified episodic memory, where verification accuracy of 85\% means 15\% of stored episodes have incorrect labels (false positives but passed verification). The data purity formula quantifies the expected corruption levels and the noise tolerant learning literature proves that training on 90-95\% pure data can improve the model instead of degrading it. This grounding distinguishes CAEM from naive approaches that might assume perfect verification is always necessary for self improvement. 

When Zelikman et al. \cite{NEURIPS2022_639a9a17} introduced STaR (Self Taught Reasoner) at NeurIPS 2022, they showed successful bootstrapping using iterative self improvement. STaR loop generates rationales for problems, filters for correct answers, fine tunes on correct rationale-answer pairs and then repeats. A key innovation is rationalization. When the model produces a correct answer without valid reasoning, it generates a rationale conditioned on the correct answer. This data augmentation stops the model from learning problems that it already knows how to solve.

In CommonsenseQA test, STaR achieved +35.9\% improvement over few-shot baseline through iterative refinement. The method shows that bootstrapping from model's own output works when coupled with appropriate filtering. The direct parallel to CAEM is evident because it generate predictions, verify correctness, fine tune on verified examples and repeat. CAEM extends STaR by adding comprehensive multi-layer verification rather than simple correctness filtering and by storing verified examples in episodic memory for future retrieval instead of solely using them for training.

The recent work of examining intrinsic self-correction capabilities reveals important limitations. Huang et al. \cite{huang-etal-2023-large} found at EMNLP 2023 that prompting LLMs to review and correct their own responses rarely improves accuracy if we do not use external feedback. Kamoi et al. \cite{kamoi-etal-2024-llms} analyzed when LLMs can correct mistakes at TACL 2024. They found that intrinsic self-correction often fails or degrades performance. The bottleneck is reliable feedback generation. So, models that produce incorrect initial answers usually cannot generate reliable correction signals without external validation. 

This limitation motivates CAEM's architectural approach. Instead of depending on intrinsic self correction using prompting alone, CAEM employs external verification through NLI models, consistency checking and semantic entropy analysis. These external signals provide the reliable feedback needed for effective self improvement. The multi-layer verification achieves 85-90\% accuracy which is substantially more trustable than model's self evaluation. This enables the virtuous improvement cycle where better generation leads to better training data and it leads to even better generation. 

The self-improvement literature also points out successful cases. When external feedback is available (for example code execution results, formal verification, retrieval of sources), models can effectively adapt the corrections. CAEM's Tier 3 routing with RAG provides exactly the same external feedback by retrieving evidence that grounds generation and enables meaningful verification. The system learns not just from its own outputs but from the interaction between generation and evidence based verification.

Domain adaptation research provides additional relevant findings. Gururangan et al. \cite{gururangan-etal-2020-dont} demonstrated at ACL 2020 that continued pre-training on domain-specific unlabeled data (domain-adaptive pretraining) followed by task-specific fine-tuning outperforms direct fine-tuning on limited labeled data. While CAEM does not perform domain-adaptive pretraining due to computational constraints, the finding supports iterative refinement: progressive training on accumulated verified examples enables gradual domain adaptation.

The distinction between CAEM's approach and standard continual learning deserves emphasis. Traditional continual learning addresses learning sequential tasks A, B, C without forgetting previous tasks. CAEM's scenario differs: the task remains constant (question answering), but the model progressively improves at that task through accumulated verified examples. This setting is closer to curriculum learning and iterative refinement than task-sequential continual learning, though catastrophic forgetting concerns still apply since fine-tuning on domain-specific verified examples risks degrading general capabilities.

Synthesis of these methods informs CAEM's self-improvement loop design. EWC and L2 regularization prevent catastrophic forgetting of general capabilities. Noisy label learning provides theoretical justification for training on imperfectly verified data. STaR demonstrates successful bootstrapping patterns. Self-correction research validates the need for external verification rather than intrinsic review. Together, these insights enable CAEM's three-cycle improvement achieving 60% to 87% accuracy progression while maintaining over 95\% general capability retention, addressing the dual requirements of task improvement and capability preservation.

\subsection{Evaluation Benchmarks}

Rigorous evaluation requires benchmarks testing diverse reasoning capabilities with objective correctness criteria. This subsection reviews four established datasets selected for CAEM evaluation, explaining how each tests different aspects of hallucination reduction.

Yang et al. \cite{yang-etal-2018-hotpotqa} introduced HotpotQA at EMNLP 2018, providing 113,000 question-answer pairs requiring reasoning across multiple Wikipedia documents. Questions are specifically designed to require multi-hop reasoning: answering requires connecting information from at least two documents that individually provide insufficient information. For example, "When was the director of Inception born?" requires first identifying Christopher Nolan as the director, then finding his birthdate. Each question includes sentence-level supporting fact annotations indicating which sentences from retrieved documents are necessary for answering.

The benchmark provides two settings: distractor setting includes nine distractor documents alongside the two gold documents, requiring both retrieval and reasoning; fullwiki setting requires searching all of Wikipedia. Human performance reaches 91.40\% F1 in the distractor setting, while best models achieve approximately 59\%, establishing substantial headroom for improvement. For CAEM, HotpotQA tests whether verified episodic memory and multi-hop reasoning chains improve accuracy on complex reasoning requiring evidence integration. The supporting fact annotations enable analysis of whether the system correctly identifies relevant evidence during verification.

Geva et al. \cite{10.1162/tacl_a_00370} introduced StrategyQA at TACL 2021, presenting 2,780 questions requiring implicit reasoning where decomposition steps are not mentioned in questions. For example, "Could a llama birth twice during the War in Vietnam?" requires knowing Vietnam War duration (implicit: 1955-1975) and llama gestation period (implicit: approximately 11 months), then computing whether two births fit the timeframe (implicit: 22 months needed, 20 years available, therefore yes). Human accuracy reaches 87\% while best models achieve approximately 66\%.

The implicit reasoning requirement tests whether models genuinely understand underlying concepts or merely perform shallow pattern matching. For CAEM, StrategyQA evaluates self-consistency verification effectiveness, as implicit reasoning should produce consistent conclusions across multiple reasoning paths if understanding is genuine. The benchmark also tests whether iterative fine-tuning improves world knowledge application, as questions require retrieving and applying multiple facts not explicitly mentioned.

Thorne et al. \cite{thorne-etal-2018-fever} created the FEVER dataset at NAACL 2018 for fact verification, providing 185,445 claims generated from Wikipedia requiring classification as SUPPORTS, REFUTES, or NOT ENOUGH INFO relative to Wikipedia evidence. Claims are specifically designed to test factual knowledge, such as "Barack Obama was the 44th President of the United States" (SUPPORTS) or "The Beatles formed in Manchester" (REFUTES, they formed in Liverpool). The NOT ENOUGH INFO category proves particularly valuable, testing whether systems can recognize insufficient evidence rather than hallucinating conclusions.

FEVER directly evaluates NLI-based verification accuracy, as the task structure matches premise-hypothesis entailment checking. Human accuracy reaches approximately 92\%, while best systems achieve 89\% in the shared task. For CAEM, FEVER tests whether NLI verification Layer 1 achieves target 85-88\% accuracy and whether the system appropriately handles the NOT ENOUGH INFO category by expressing uncertainty rather than forcing conclusions.

Lin et al. \cite{lin-etal-2022-truthfulqa} designed TruthfulQA specifically to measure how models mimic human falsehoods, providing 817 questions across 38 categories in health, law, finance, politics, and conspiracies. Questions are adversarially constructed to elicit common misconceptions: "What happens if you crack your knuckles?" where the false but common answer is "you get arthritis" and the true answer is "nothing harmful happens; no proven link to arthritis." Best-answer evaluation measures whether models generate truthful responses, while true-false evaluation tests whether models rate true statements as more likely than false ones.

The inverse scaling finding (larger models often less truthful) makes TruthfulQA the most challenging benchmark for parametric scaling approaches. Human performance reaches 94\% while GPT-3 175B achieves only 21-25\% on the truthful-and-informative metric (58\% on truthfulness alone), demonstrating substantial gaps. For CAEM, TruthfulQA critically tests semantic entropy verification, as misconceptions often produce high semantic entropy (models generate both correct and false answers with similar confidence). The benchmark validates whether multi-layer verification catches systematic confabulation better than single-method approaches.

Benchmark complementarity ensures comprehensive evaluation. HotpotQA tests multi-hop reasoning with explicit evidence. StrategyQA tests implicit reasoning with world knowledge. FEVER tests direct fact verification and uncertainty expression. TruthfulQA tests resistance to misconceptions and confabulation. Together, they cover the spectrum of factual reasoning tasks where hallucination reduction matters most. Success across all four benchmarks demonstrates that CAEM's improvements generalize across reasoning types rather than overfitting to specific task characteristics.

\section{Summary of Key Findings}

The literature review not only indicates the necessity but also provides technical premises for the architectural orientation of CAEM to reduce hallucinations. The studies of hallucination measurement show a prevalence between upto 90\% with all domains. TruthfulQA's inverse scaling idea \cite{lin-etal-2022-truthfulqa} where GPT-3 175B achieved only 21-25\% on the combined truthful and informative metric showed parametric scaling alone is insufficient for truthfulness. FActScore \cite{min-etal-2023-factscore} implies that even competent models give 58\% supported atomic facts. It supports the need of granular verification. These findings promote the multi-component architecture intervention of CAEM. The verification literature provides the literature of complementary strengths. The interpretable reasoning traces of the chain of thought \cite{NEURIPS2022_9d560961}, semantic entropy \cite{farquhar2024semantic} with 0.790 AUROC of entailment detection with meaning level clustering, and NLI verification \cite{liu2019roberta} have 90.8\% of entailment classification. CAEM combines them by three layer checking with a result of an anticipated 85-90\% accuracy. According to research conducted on confidence estimation, single signals have a maximized 0.5-0.6 AUROC \cite{xiong2023can} which is barely even above a random value. Systematic overconfidence \cite{pmlr-v70-guo17a} on the other hand must have post-hoc calibration by the temperature scaling. The RLHF training \cite{kadavath2022language} compromises the calibration and this shows that the models, which are instruction tuned such as Flan T5, give reliable signals. CAEM is a combination of four complementary signal of confidence (token probability, MC dropout, self-consistency, semantic entropy) which uses white box, gray box and black box approaches where temperature scaling is used to calibrate the routing decisions.

Memory augmented architectures determine the existence of extrinsic knowledge storage. It is proven that structured reasoning is possible with differentiable memory in Neural Turing Machines \cite{graves2014neural}, Memory Networks \cite{weston2015memory}, and DNCs \cite{graves2016differentiable}. RAG \cite{NEURIPS2020_6b493230}, however, demonstrates that the non-parametric knowledge is more accurate (44.5\% vs 41.5\% on natural questions) and does not require retraining to update the knowledge. FAISS \cite{8733051} and Sentence-BERT \cite{reimers-gurevych-2019-sentence} allow to search similarities efficiently that it is possible to compute 10,000 sentences in 5 seconds. The CAEM expands RAG, where routing and similarity are maintained using a score based on confidence, and reasoning chains are stored after verification to achieve progressively higher efficiency. The continual learning studies give theoretical support for the iterative improvement. The EWC \cite{doi:10.1073/pnas.1611835114} maintains 90\% performance among sequential tasks. Noisy label learning \cite{natarajan2013learning} achieves 88\% accuracy with 40\% label corruption. There is an improvement in the STaR \cite{NEURIPS2022_639a9a17} of +35.9\% with bootstrapping. According to self correction research \cite{huang-etal-2023-large, kamoi-etal-2024-llms}, intrinsic correction cannot work without external feedback which CAEM fills in with multi-layer external verification. The CAEM uses regularization of L2 to mitigate catastrophic forgetting, as well as iterative improvement through cycles. The four benchmarks that we selected have the capability of providing coverage of QA reasoning. HotpotQA \cite{yang-etal-2018-hotpotqa} tests multi-hop reasoning (91.4\% human vs 59\% model F1), StrategyQA \cite{10.1162/tacl_a_00370} tests implicit reasoning using world knowledge, Fever \cite{thorne-etal-2018-fever} tests fact verification and expressing uncertainty and lastly, TruthfulQA \cite{lin-etal-2022-truthfulqa} tests resisting misconceptions. Success across all benchmarks will generalize rather than task specific overfitting with primary metrics (exact match, F1, hallucination rate) and secondary metrics (verification accuracy, data purity, latency, memory usage) that will provide a comprehensive evaluation. 

The literature establishes that hallucination remains unsolved despite massive investment, parametric scaling proves insufficient, and individual methods achieve limited success. CAEM addresses these limitations through verified episodic memory, multi-layer verification, multi-signal confidence estimation, adaptive routing, and iterative self-improvement, grounded in both empirical validation plans and theoretical analysis of data purity and convergence dynamics.